% Generated draft: requires mathematical and source review.
\documentclass[11pt]{article}
\usepackage[margin=25mm]{geometry}
\usepackage{fontspec}
\setmainfont{Latin Modern Roman}
\usepackage{amsmath,amssymb,mathtools,unicode-math}
\setmathfont{Latin Modern Math}
\usepackage{enumitem,needspace,xurl,hyperref}
\setlength{\parindent}{0pt}
\setlength{\parskip}{5pt}
\setlength{\emergencystretch}{3em}
\newcommand{\N}{\mathbb{N}}
\newcommand{\Z}{\mathbb{Z}}
\newcommand{\Q}{\mathbb{Q}}
\newcommand{\R}{\mathbb{R}}
\newcommand{\C}{\mathbb{C}}
\newcommand{\F}{\mathbb{F}}
\newcommand{\A}{\mathbb{A}}
\newcommand{\PP}{\mathbb{P}}
\newcommand{\E}{\mathbb{E}}
\newcommand{\eps}{\varepsilon}
\newcommand{\dd}{\,\mathrm{d}}
\begin{document}
\section*{UnsolvedMath ID: 30004718}
\textbf{Source problem code:} catalogue entry
\subsection{Definitions and mathematical statement}
Let \(n\) be an integer for which the displayed expressions are defined; in particular, they require \(n>1\). Let \(Y\) be a closed smooth \((n-1)\)-manifold, where closed means compact without boundary, and suppose that
\[
n-1\ne 4,
\qquad
Y\text{ admits no Riemannian metric of everywhere positive scalar curvature}.
\]
Set
\[
X=Y\times[-1,1],\qquad
\partial_-X=Y\times\{-1\},\qquad
\partial_+X=Y\times\{1\},\qquad
c_n=\sqrt{\frac{n-1}{n}}.
\]
For a smooth Riemannian metric \(g\) on \(X\), let \(\operatorname{Sc}_g\) denote its scalar curvature and define
\[
\operatorname{width}(X,g)
=\inf\bigl\{L_g(\gamma):\gamma\text{ is a piecewise smooth curve in }X,
\ \gamma(0)\in\partial_-X,\ \gamma(1)\in\partial_+X\bigr\}.
\]
Write \(H_-\) and \(H_+\) for the mean curvatures of the respective boundary components. The supplied record does not specify the normal directions, sign convention, or whether mean curvature is the trace or the average of the principal curvatures. The boundary inequalities below therefore retain the source's unresolved convention and are to hold pointwise.

\textbf{Reported conjecture, with unresolved qualifications.} For every such \(Y\) and \(g\), every real \(\sigma>0\), and every pair of real parameters \(\ell_-\), \(\ell_+\) satisfying
\[
-\frac{\pi c_n}{\sqrt{\sigma}}
<\ell_-<\ell_+
<\frac{\pi c_n}{\sqrt{\sigma}},
\]
suppose that
\[
\begin{aligned}
\operatorname{Sc}_g(x)&\ge\sigma &&\text{for every }x\in X,\\
H_-(p)&\ge\sqrt{\sigma}\,c_n
\tan\!\left(\frac{\sqrt{\sigma}\,\ell_-}{2c_n}\right)
&&\text{for every }p\in\partial_-X,\\
H_+(p)&\ge-\sqrt{\sigma}\,c_n
\tan\!\left(\frac{\sqrt{\sigma}\,\ell_+}{2c_n}\right)
&&\text{for every }p\in\partial_+X,\\
\operatorname{width}(X,g)&\ge\ell_+-\ell_-.
\end{aligned}
\]
Must there exist a Riemannian metric \(g_Y\) on \(Y\) with \(\operatorname{Ric}_{g_Y}=0\) and a Riemannian isometry
\[
(X,g)\cong
\left(Y\times[\ell_-,\ell_+],\ \varphi(t)^2g_Y+\dd t^2\right),
\qquad
\varphi(t)=
\cos\!\left(\frac{\sqrt{\sigma}\,t}{2c_n}\right)^{n/2}?
\]
The parameter bounds make the cosine positive throughout the target interval. The exponent \(n/2\) is reproduced exactly as supplied; its compatibility with the intended rigidity model and the other constants requires source verification. No spin, index, orientability, or connectedness assumption is added.

The supplied expanded statement asks this implication under a width lower bound. Its shorter original formulation instead says that ``equality forces'' the displayed model without specifying the equality concerned. The expanded implication is recorded here as the supplied target, with that discrepancy unresolved. The record also gives no explicit dimension range beyond \(n-1\ne4\); the restriction \(n>1\) above only expresses the requirement that its formulas be defined.
\subsection{Short English statement}
For a band over a closed manifold admitting no positive scalar curvature metric, the supplied conjecture asks whether the stated scalar curvature, boundary mean curvature, and width bounds force a warped product with a Ricci-flat base. The boundary mean curvature convention, the intended equality formulation, and the supplied warping exponent require clarification. This is a statement-only draft requiring human mathematical review.
\subsection{Sources}
\begin{itemize}
\item Local UnsolvedMath registry: supplied complete record with permanent ID \(30004718\), source problem code \texttt{catalogue entry}, title \emph{Rigidity in Scalar-Curvature Band Width Inequalities}, in the Oberwolfach Reports Open Problems collection. The record gives proposed year \(2021\) and names no individual proposer.
\item Supplied report citation: \emph{Analysis, Geometry and Topology of Positive Scalar Curvature Metrics} (2022), Oberwolfach Reports, pp.~1579--1649, DOI \texttt{10.4171/owr/2021/30}: \url{https://doi.org/10.4171/owr/2021/30}. The supplied verification note identifies the extracted statement as Conjecture~4 and reports that the scalar curvature, boundary mean curvature, parameter-range, and width hypotheses were restored after inspection. This is a supplied verification claim, not an inspection performed for this draft.
\item Supplied literature source, labeled ``Oberwolfach Reports 2022'': \url{https://ems.press/content/serial-article-files/46909}.
\item Supplied literature source, labeled ``Rade, doctoral thesis on band rigidity'': \url{https://opus.bibliothek.uni-augsburg.de/opus4/frontdoor/index/index/year/2023/docId/99573}.
\item Supplied literature assessment, checked on \texttt{2026-08-22}: the record labels the problem partially solved and states that rigidity is proved under spin, nonzero index, and related model hypotheses, while the conjectured rigidity statement in the extracted generality remains open. This is supplied dated curation derived from the OWR record, not an independently verified literature assessment. Its time-sensitive status and cited claims require independent verification.
\end{itemize}
\end{document}
